% !TEX root = paper.tex
% =============================================================
% RAR -- Related Work v0
% Source-of-truth: project_context.md (binding constraints).
% Stage 3 of paper-writing skill. Draft 0 intro frozen 2026-07-09.
% Last updated: 2026-07-09
%
% Three categories (locked):
%   (1) CIL methods -- LwF, iCaRL, BiC, DER family
%   (2) Traffic-classification backbones -- ET-BERT, AppScanner, etc.
%   (3) MEMENTO -- mentioned briefly as complementary work, not as a baseline
%
% Conventions:
%   - KDD venue: integrated related work, colon-style subtitles,
%     no \smartparagraph, no em-dashes.
%   - Move 1 (Category Clustering): 2-4 categories by approach type.
%   - Move 2 (Per-Category Limitation): STRUCTURAL, not quantitative.
%   - Move 3 (Positioning Sentence): preempt "different from MEMENTO?"
% =============================================================

\section{Related Work}\label{sec:related}

% Section opener: positions RAR in the CIL design space.

\subsection{Class-Incremental Learning}

% Topic sentence: cluster and limit
Classical CIL methods commonly combine distillation, calibration, and replay. The representative methods considered here generally use class-level or dataset-level balancing rather than an input-dependent old/new routing score.

% Distillation-based (LwF, iCaRL)
Distillation-based methods (LwF~\cite{li2017lwf}, iCaRL~\cite{rebuffi2017icarl}) match the soft outputs of the old model using a globally weighted objective. The structural limitation is that a fixed coefficient does not adapt the preservation pressure to each sample's compatibility with historical classes.

% Calibration-based (BiC)
Calibration-based methods such as BiC~\cite{wu2019bic} learn a linear correction for the new-class logits after incremental training. This correction is shared across inputs and therefore addresses systematic old/new bias rather than input-dependent routing.

% Replay-based (DER)
Replay-based methods (DER~\cite{yan2021der}) store old samples or features and replay them during incremental training. The structural limitation is the buffer: a fixed buffer cannot represent the full old-class distribution, and the model's behavior on unstored samples inherits the buffer's under-representation.

% What RAR does
RAR adds a per-sample old-versus-new score derived from reconstruction error. The score gates the two prediction paths and weights the distillation and margin losses, while replay supplies the old samples used to fit the reconstructor and train the incremental model.

\subsection{Task Inference and Expert Routing}

Input-dependent routing predates RAR. Expert Gate~\cite{aljundi2017expertgate} trains an autoencoder for each sequentially learned task and selects the expert whose autoencoder yields the lowest reconstruction error for a test sample. It establishes reconstruction error as a task-inference signal. A complementary line connects CIL to task prediction and OOD detection: CLOM~\cite{kim2022clom} trains task-specific OOD models, a theoretical analysis formalizes within-task and task-ID prediction as the two components of CIL~\cite{kim2022theoretical}, MORE~\cite{kim2022more} builds task-specific heads with OOD replay, and TPL~\cite{lin2024tpl} improves task prediction through likelihood ratios.

RAR does not claim reconstruction-based task inference or multi-head CIL as new. It addresses encrypted-traffic CIL with a different integration: a single reconstructor fitted to replayed historical traffic produces a continuous old/new compatibility score; the score weights preservation and acquisition losses and softly gates two disjoint class groups; calibration then supports joint prediction over all seen traffic classes. The contribution is this reconstruction-conditioned training and inference framework, rather than any component in isolation.

\subsection{Traffic Classification Backbones}

% Topic sentence: cluster and limit
Encrypted-traffic classification research has produced a series of strong backbones (ET-BERT~\cite{lin2022etbert}, AppScanner, ETC-PS, FlowLens, CNN-based, CNN-emb, EBSNN, FS-Net, GraphDApp, MM4flow) that pre-train on packet-level features and reach high accuracy on static label sets. The structural limitation of this line of work is that it solves representation learning, not incremental learning: a stronger backbone still forgets old classes when its classification head expands to accommodate new ones. We quantify this in \S\ref{sec:eval-bgroup}: naive incremental training forgets across all nine backbones on NUDT-Mobile, and RAR (on top of ET-BERT) breaks out of the band.

% What RAR does
RAR's reconstruction-error routing operates on frozen encoder features and therefore depends on their separability. The routing mechanism does not require the encoder itself to be incremental-aware, and its per-sample score can be incorporated into frozen-encoder CIL pipelines. Recent work on incremental representation learning, such as MEMENTO~\cite{memento2024}, instead adapts the backbone to the class stream. RAR studies the complementary routing problem; combining the two forms of adaptation remains future work.
